\chapter{การประมวลผลทางวิทยาศาสตร์และการควบคุมสะเต็มด้วย Python และ PyOpenXR}
\label{chap:ch05}

\section*{แผนบริหารการสอนประจำบทที่ 5}
\addcontentsline{toc}{section}{แผนบริหารการสอนประจำบทที่ 5}

\begin{planbox}[colframe=darkNavy, colbacktitle=darkNavy]{1. หัวข้อเนื้อหาประจำบท}
\begin{enumerate}[topsep=2pt, itemsep=1pt, partopsep=0pt, parsep=0pt, leftmargin=1.5em]
    \item สถาปัตยกรรมตัวเชื่อมต่อ PyOpenXR และการจัดการหน่วยความจำผ่าน ctypes
    \item การประมวลผลข้อมูลทางวิทยาศาสตร์สมรรถนะสูงด้วย NumPy และ Zero-Copy Buffers
    \item การแยกเธรดการคำนวณฟิสิกส์ออกจากเธรดการเรนเดอร์ภาพ (Multithreaded Architecture)
    \item การพล็อตข้อมูลและการแสดงผลสนามเวกเตอร์ทางฟิสิกส์แบบเรียลไทม์ในโลกเสมือน
\end{enumerate}
\end{planbox}

\begin{planbox}[colframe=openxrTeal, colbacktitle=openxrTeal]{2. วัตถุประสงค์เชิงพฤติกรรม}
\begin{enumerate}[topsep=2pt, itemsep=1pt, partopsep=0pt, parsep=0pt, leftmargin=1.5em]
    \item อธิบายกลไกการเชื่อมต่อระหว่าง Python C-API กับ OpenXR Native Loader ได้
    \item จัดการอาร์เรย์พิกัดสามมิติด้วย NumPy โดยไม่ต้องคัดลอกข้อมูลซ้ำซ้อนได้
    \item ออกแบบโครงสร้างมัลติเธรดเพื่อรักษาอัตราเฟรมภาพที่ 90 FPS ได้อย่างเสถียร
    \item สร้างระบบแสดงผลเส้นแรงแม่เหล็กและปริมาณฟิสิกส์สามมิติเชิงโต้ตอบได้
\end{enumerate}
\end{planbox}

\newpage

\begin{planbox}[colframe=spatialBlue, colbacktitle=spatialBlue]{3. วิธีสอนและกิจกรรมการเรียนการสอน}
\begin{enumerate}[topsep=2pt, itemsep=1pt, partopsep=0pt, parsep=0pt, leftmargin=1.5em]
    \item การบรรยายโครงสร้าง PyOpenXR พร้อมการเขียนโค้ดสด (Live Coding Demonstration)
    \item กิจกรรมการคำนวณสนามไฟฟ้าและแม่เหล็กรอบขดลวดโซเลนอยด์ด้วย NumPy
    \item การฝึกออกแบบ Ring Buffer เพื่อถ่ายโอนข้อมูลระหว่างเธรดฟิสิกส์และเธรดเรนเดอร์
    \item การอภิปรายปัญหาและแนวทางแก้ไข Global Interpreter Lock (GIL) ในภาษา Python
\end{enumerate}
\end{planbox}

\begin{planbox}[colframe=rbruGold!90!black, colbacktitle=rbruGold!90!black]{4. สื่อการเรียนการสอน}
\begin{enumerate}[topsep=2pt, itemsep=1pt, partopsep=0pt, parsep=0pt, leftmargin=1.5em]
    \item สไลด์บรรยายอิเล็กทรอนิกส์และเอกสารโค้ดมาตรฐานประจำบทที่ 5
    \item สภาพแวดล้อมเสมือน Python (Virtualenv) ที่ติดตั้ง PyOpenXR และ OpenGL
    \item ชุดข้อมูลการจำลองสนามเวกเตอร์และวิถีการเคลื่อนที่ของอนุภาค
    \item จอแสดงผลข้อมูลทางวิทยาศาสตร์แบบสามมิติ (3D Spatial Telemetry HUD)
\end{enumerate}
\end{planbox}

\begin{planbox}[colframe=darkSlate, colbacktitle=darkSlate]{5. การวัดและประเมินผล}
\begin{enumerate}[topsep=2pt, itemsep=1pt, partopsep=0pt, parsep=0pt, leftmargin=1.5em]
    \item การตรวจประเมินโปรแกรมคำนวณและแสดงผลสนามแม่เหล็กสามมิติ
    \item การวัดประสิทธิภาพอัตราเฟรมและความหน่วงของระบบมัลติเธรดที่พัฒนา
    \item การทดสอบข้อเขียนเชิงสถาปัตยกรรมซอฟต์แวร์ทางวิทยาศาสตร์
\end{enumerate}
\end{planbox}
\clearpage

\section{สถาปัตยกรรม PyOpenXR และการเชื่อมต่อระดับลึก}

ภาษา Python เป็นภาษาหลักในการคำนวณทางวิทยาศาสตร์และปัญญาประดิษฐ์ การผสานพลังของ Python เข้ากับมาตรฐาน OpenXR ผ่านแพ็กเกจ \texttt{pyopenxr} ช่วยให้นักวิจัยและนักการศึกษาสามารถนำโมเดลคณิตศาสตร์ที่ซับซ้อนมาแสดงผลเชิงพื้นที่ได้อย่างรวดเร็ว

\begin{figure}[htbp]
\centering
\begin{tikzpicture}[node distance=1.6cm, auto]
    % Boxes
    \node (py) [rectangle, rounded corners, draw=spatialBlue, fill=spatialBlue!15, line width=1.2pt, minimum width=10cm, minimum height=1.1cm] {\textbf{Python Application Layer} (Scientific Simulation, NumPy, SciPy)};
    \node (ctypes) [rectangle, rounded corners, draw=openxrTeal, fill=openxrTeal!15, line width=1.2pt, minimum width=10cm, minimum height=1.1cm, below of=py] {\textbf{PyOpenXR Binding Layer} (Ctypes FFI / Zero-Copy Memory Views)};
    \node (native) [rectangle, rounded corners, draw=darkNavy, fill=darkNavy!15, line width=1.2pt, minimum width=10cm, minimum height=1.1cm, below of=ctypes] {\textbf{Native OpenXR Loader} (\texttt{libopenxr\_loader.so} / \texttt{.dll})};
    \node (gpu) [rectangle, rounded corners, draw=amberGlow, fill=amberGlow!15, line width=1.2pt, minimum width=10cm, minimum height=1.1cm, below of=native] {\textbf{Hardware Graphics \& XR Runtime} (OpenGL / Vulkan Swapchains)};

    % Arrows
    \draw [->, line width=1.3pt, spatialBlue] (py) -- (ctypes) node[midway, right, font=\footnotesize]{Direct Pointer Access};
    \draw [->, line width=1.3pt, openxrTeal] (ctypes) -- (native) node[midway, right, font=\footnotesize]{C-ABI Function Calls};
    \draw [->, line width=1.3pt, darkNavy] (native) -- (gpu) node[midway, right, font=\footnotesize]{Stereoscopic Frame Submission};
\end{tikzpicture}
\caption{สถาปัตยกรรมการเชื่อมต่อของ PyOpenXR ผ่าน ctypes FFI สู่รันไทม์ระดับฮาร์ดแวร์}
\label{fig:pyopenxr_pipeline}
\end{figure}

\section{การคำนวณสมรรถนะสูงและการจัดการบัฟเฟอร์}

ความท้าทายหลักของการใช้ภาษา Python คือความเร็วในการวนลูป ดังนั้น การคำนวณตำแหน่งอนุภาคหรือสนามเวกเตอร์ต้องกระทำผ่านเวกเตอร์ไรเซชัน (Vectorization) ของ NumPy และส่งผ่านข้อมูลไปยัง GPU บัฟเฟอร์โดยตรงโดยไม่ผ่านการคัดลอก (Zero-Copy Transfer)

\begin{table}[htbp]
\centering
\caption{การเปรียบเทียบเวลาในการประมวลผลสนามเวกเตอร์ 10,000 จุดในภาษา Python}
\label{tab:python_performance}
\begin{tabularx}{\textwidth}{lYY}
\toprule
\textbf{วิธีการเขียนโปรแกรม} & \textbf{เวลาในการประมวลผล (ms)} & \textbf{อัตราเฟรมสูงสุดที่ทำได้ (FPS)} \\
\midrule
Python Standard For-loop & 145.2 ms & $\sim 6.8\text{ FPS}$ (กระตุกรุนแรง) \\
NumPy Vectorized Array & 2.1 ms & $> 90\text{ FPS}$ (ลื่นไหลระดับสากล) \\
NumPy + Numba JIT Compilation & 0.6 ms & $> 120\text{ FPS}$ (สมรรถนะสูงสุด) \\
\bottomrule
\end{tabularx}
\end{table}

\begin{workedexamplebox}{การคำนวณสนามแม่เหล็กแบบเวกเตอร์รอบเส้นลวดนำกระแสไฟฟ้า}
เส้นลวดยาวตรงวางตัวอยู่ตามแนวแกน $Z$ และมีกระแสไฟฟ้า $I = 10\text{ A}$ ไหลในทิศทาง $+Z$ จงเขียนฟังก์ชันเวกเตอร์คำนวณเวกเตอร์สนามแม่เหล็ก $\mathbf{B}$ ที่จุดพิกัด $\mathbf{r} = (x, y, 0)$ และคำนวณค่าสนามแม่เหล็กที่จุด $(0.05, 0.05, 0)\text{ m}$

\textbf{วิธีทำ:}
ตามกฎของบิโอต์-ซาวารต์และกฎของแอมแปร์ สนามแม่เหล็กจากเส้นลวดตรงยาวมีขนาด
\begin{equation}
B = \frac{\mu_0 I}{2\pi r}
\end{equation}
โดยที่ $\mu_0 = 4\pi \times 10^{-7}\text{ T}\cdot\text{m/A}$ และ $r = \sqrt{x^2 + y^2}$ เวกเตอร์หนึ่งหน่วยในแนวสัมผัสคือ $\hat{\boldsymbol{\phi}} = \left(-\frac{y}{r}, \frac{x}{r}, 0\right)$
ดังนั้น เวกเตอร์สนามแม่เหล็กในพิกัดคาร์ทีเซียนคำนวณได้จาก
\begin{equation}
\mathbf{B}(x, y) = \frac{\mu_0 I}{2\pi (x^2 + y^2)} (-y\hat{\mathbf{i}} + x\hat{\mathbf{j}})
\end{equation}
ที่จุด $(0.05, 0.05, 0)\text{ m}$:
\begin{align}
r^2 &= (0.05)^2 + (0.05)^2 = 0.0025 + 0.0025 = 0.0050\text{ m}^2 \\
B_{\text{scale}} &= \frac{(4\pi \times 10^{-7})(10)}{2\pi (0.0050)} = \frac{2 \times 10^{-6}}{0.0050} = 4.0 \times 10^{-4}\text{ T}
\end{align}
จะได้เวกเตอร์สนามแม่เหล็ก
\begin{equation}
\mathbf{B} = 4.0 \times 10^{-4} (-0.05\hat{\mathbf{i}} + 0.05\hat{\mathbf{j}}) / \sqrt{0.0050} \approx (-2.83\times 10^{-5}\hat{\mathbf{i}} + 2.83\times 10^{-5}\hat{\mathbf{j}})\text{ T}
\end{equation}
ข้อมูลเวกเตอร์นี้สามารถถูกเรนเดอร์เป็นลูกศร 3 มิติในโลกเสมือนจริงของ OpenXR เพื่อให้ผู้เรียนเห็นทิศทางการวนของสนามแม่เหล็กได้อย่างชัดเจน
\end{workedexamplebox}

\section{สถาปัตยกรรมแบบแยกเธรดการประมวลผล}

เพื่อให้การเรนเดอร์ภาพผ่าน OpenXR ไม่ถูกรบกวนจากการคำนวณทางวิทยาศาสตร์ที่มีความซับซ้อน เราต้องแยกวงรอบการทำงานออกเป็น 2 เธรดหลัก ได้แก่
\begin{enumerate}
    \item \textbf{Physics Simulation Thread:} คำนวณสมการอนุพันธ์ การเคลื่อนที่ของอนุภาค และการตรวจจับการชน โดยทำงานด้วยความถี่คงที่ (Fixed Timestep เช่น $\Delta t = 0.01\text{ s}$)
    \item \textbf{XR Rendering Thread:} ทำหน้าที่รับข้อมูลพิกัดล่าสุดผ่านโครงสร้างหน่วยความจำแบบล็อกอิสระ (Lock-Free Shared Memory) และส่งเฟรมภาพไปยังอุปกรณ์แสดงผลด้วยความถี่ของฮาร์ดแวร์ ($90\text{ Hz}$)
\end{enumerate}

\begin{aipromptbox}[การสร้างท่อส่งข้อมูลวิทยาศาสตร์ NumPy สู่ GPU Buffer แบบ Zero-Copy ด้วย AI]
\textbf{วัตถุประสงค์:} ออกแบบท่อส่งข้อมูล (Data Pipeline) ในภาษา Python ที่คำนวณสนามเวกเตอร์ 3 มิติด้วย NumPy และส่งต่อไปยัง OpenGL/Vulkan Vertex Buffer Object (VBO) โดยตรงผ่านตัวชี้หน่วยความจำ (Zero-Copy Transfer)\\
\textbf{โครงสร้าง CLEAR Prompt:}
\begin{itemize}[topsep=1pt, itemsep=1pt, leftmargin=1.2em]
    \item \textbf{Context (บริบท):} การแสดงผลเส้นแรงแม่เหล็กและอนุภาค 50,000 จุดบน OpenXR แบบเรียลไทม์ที่ 90 FPS ในภาษา Python
    \item \textbf{Logic (ตรรกะ):} ใช้ \texttt{numpy.ndarray.ctypes.data} เพื่อเข้าถึงที่อยู่หน่วยความจำจริง และผูกเข้ากับ \texttt{glBufferSubData} โดยไม่ผ่านการแปลงชนิดข้อมูลซ้ำ
    \item \textbf{Expectation (ผลลัพธ์ที่ต้องการ):} ฟังก์ชัน Python \texttt{update\_vector\_field\_vbo(vbo\_id, positions, vectors)} พร้อมตัวอย่างการวัดเวลาดำเนินการระดับไมโครวินาที
    \item \textbf{Action \& Restriction (ข้อกำหนด):} หน่วยความจำต้องเรียงลำดับแบบ C-contiguous (\texttt{order='C'}) และข้อมูลต้องเป็นชนิด \texttt{np.float32} เท่านั้น
\end{itemize}
\end{aipromptbox}

\begin{codebox}[การจำลองสนามเวกเตอร์ความเร็วสูงด้วย NumPy ใน PyOpenXR]
import numpy as np

def compute_magnetic_field_grid(current, grid_x, grid_y):
    # grid_x, grid_y เป็น 2D NumPy Arrays
    mu_0 = 4.0 * np.pi * 1e-7
    r_squared = grid_x**2 + grid_y**2
    # ป้องกันการหารด้วยศูนย์ที่จุดแกนลวด
    r_squared = np.maximum(r_squared, 1e-6)

    factor = (mu_0 * current) / (2.0 * np.pi * r_squared)
    Bx = -factor * grid_y
    By = factor * grid_x
    Bz = np.zeros_like(Bx)

    # ส่งคืนอาร์เรย์สามมิติสำหรับเรนเดอร์ใน OpenXR
    return np.stack([Bx, By, Bz], axis=-1)
\end{codebox}

\section*{สรุปท้ายบทที่ 5}
\addcontentsline{toc}{section}{สรุปท้ายบทที่ 5}

การผสานพลังระหว่างภาษา Python และมาตรฐาน OpenXR เปิดขอบฟ้าใหม่ให้กับการศึกษาสะเต็มศึกษา ด้วยการใช้ความสามารถในการคำนวณเมทริกซ์และเวกเตอร์ของ NumPy ร่วมกับสถาปัตยกรรมมัลติเธรด นักพัฒนาสามารถสร้างห้องปฏิบัติการเสมือนจริงที่คำนวณสมการฟิสิกส์ได้อย่างแม่นยำ พร้อมทั้งแสดงผลข้อมูลทางวิทยาศาสตร์แบบสามมิติได้อย่างลื่นไหลและมีประสิทธิภาพสูงสุด

\section*{แบบฝึกหัดท้ายบทที่ 5}
\addcontentsline{toc}{section}{แบบฝึกหัดท้ายบทที่ 5}

\begin{enumerate}
    \item จงอธิบายบทบาทของ ctypes ในการเชื่อมต่อระหว่างฟังก์ชันของ PyOpenXR กับไลบรารีภาษา C ของ OpenXR Loader
    \item เพราะเหตุใดการเขียนโปรแกรมคำนวณตำแหน่งอนุภาค 10,000 ตัวด้วย For-loop ปกติใน Python จึงไม่สามารถใช้ในงาน OpenXR แบบเรียลไทม์ได้
    \item ออกแบบผังการทำงาน (Flowchart) ของสถาปัตยกรรมมัลติเธรดที่แยกเธรดการคำนวณฟิสิกส์ออกจากเธรดการเรนเดอร์ภาพของ OpenXR
\end{enumerate}
